\scanpage{15}\mspage{12}

\deleted{\unsure{Noter que}} $\Phi_{\unsure{0|\circ}}$ est un module libre sur $\mathbb Z$, engendré par les
\[
 \varphi_i^{\beta,\alpha}\deleted{\unsure{\pi_{i,\alpha}}}
 \qquad\begin{aligned}&0\le i\le n\\&0\le\alpha,\beta\le2n-i\end{aligned}
\]
\[
 \deleted{\left\{ L^j\pi_{i,\alpha}\quad\varphi_{i,\alpha}^{j,\beta}\ \unsure{\pi_{i,\alpha}}\quad\text{pour }
 \left\{\begin{aligned}
 &\unsure{\text{avec }}0\le i\le 2n,\ 0\le j\le2n,\ \alpha,\beta\ge0\\
 &i-2\alpha=j-2\beta\ge0\\
 &\unsure{0\le i<2}
 \end{aligned}\right.\right.}
\]
% slanted
\margin{$\mathbb Z$ peut être remplacé par n'importe quel anneau}
\deleted{$\varphi_{i,\alpha}^{j,\beta}$ est nul sur les composantes} définis par
\[
 \left|\begin{aligned}
 \varphi_i^{\beta,\alpha}&=L_0^{\beta-\alpha}\pi_{i+2\alpha,\alpha}\quad\text{si }\alpha\le\beta\\
 \varphi_i^{\beta,\alpha}&=\Lambda_0^{\alpha-\beta}\pi_{i+2\alpha,\alpha}\quad\text{si }\alpha\ge\beta
 \end{aligned}\right.
\]
\deleted{\[
 \left\{\begin{aligned}
 \varphi_{i,\alpha}^{j,\beta}&=L_0^{\beta-\alpha}\pi_{i,\alpha}=\pi_{j,\beta}L_0^{\beta-\alpha}\quad\text{si }\alpha\le\beta\\
 \varphi_{i,\alpha}^{j,\beta}&=\Lambda_0^{\alpha-\beta}\pi_{i,\alpha}=\pi_{j,\beta}\Lambda_0^{\alpha-\beta}\quad\text{si }\alpha\ge\beta
 \end{aligned}\right.
\]}
i.e. on a une base formée des
\[
 \varphi_{i'}^{\beta',\alpha'}\varphi_i^{\beta,\alpha}
 \quad\deleted{\unsure{\varphi_{i,\alpha}^{j,\beta}\varphi_{i',\alpha'}^{j',\beta'}}}
 =\left|\begin{aligned}
 &0\ \text{si }\deleted{\unsure{j\ne i'\text{ ou }2\alpha\ne\beta}}\inserted{i\ne i'\text{ ou }\beta\ne\alpha'}\\
 &\varphi_{\unsure{i|i'}}^{\beta',\alpha}\quad \text{sinon}
 \end{aligned}\right.
\]
\deleted{\[
 \left\{ L^\beta\pi_{i,\alpha}\qquad
 \left|\begin{aligned}
 &0\le i\le\unsure{2n|n},\quad\alpha\ge0\quad\text{tel que }i-2\alpha\ge0\\
 &0\le2\alpha\le i\\
 &0\le\beta
 \end{aligned}\right.\right.
\]}
% slanted
\margin{Les $\varphi_i^{\beta\alpha}$ et leurs multiplications devraient être donnés au n\textsuperscript{o} 6}

La multiplication des élémts de cette base ne dépend pas du choix de $\mathbb M_0,\unsure{\xi|\zeta}_0$ (ni du fait qu'on est sur l'anneau $\mathbb Z$ \ldots), donc l'anneau $\Phi_0$ est bien déterminé en termes de l'entier \inserted{pair} $2n$.

\uline{Proposition} Revenant à la donnée de $M^*$ comme au n\textsuperscript{o} 1, et de $\unsure{E|\mathcal E}$ ss-\deleted{algèbre}\inserted{anneau} graduée de \deleted{\unsure{$E^*$}} $\operatorname{End}(M^*)$. Soit $L\in E^{(\unsure{2|1})}$. Conditions équivalentes

(i) L'endom $L$ de $M^*$ satisfait la condition de \deleted{\unsure{Hodge}} Lefschetz, et l'opérateur $\Lambda$ associé est dans $E$. (i.e on est \unsure{sous|dans} les conditions du th. du n\textsuperscript{o} 6).

(ii) Il existe un hom d'anneaux gradués
\[
 \varphi:\Phi_0\longrightarrow E
\]
% archive sheet number: 14
